\section{The global equivalence and its central data}
\label{sec:equivalence}

We now combine the period comparison with the Plancherel identity.  The
period comparison first produces a functor out of a localization quotient.
The two Plancherel identities, one for each group, then provide its two
inverses.  We finish by describing the compatibility with localization,
components, central characters, and central torsors.

\subsection{The four localizations}

Recall the shifted levels and determinant exponents
\[
 a=-c,\qquad d=\frac1{ra},\qquad
 k_a=a-h^\vee,\qquad k_d=d-\check h^\vee,\qquad
 s=-\frac{k_a}{2h^\vee},\qquad
 s'=-\frac{k_d}{2\check h^\vee}.
\]
Thus the four bundle categories used below are
\[
 \begin{array}{ll}
 A=D(L_G^s),& A'=D(L_G^{-s}),\\
 B=D(L_{G^\vee}^{-s'}),&B'=D(L_{G^\vee}^{s'}).
 \end{array}
\]
All contain every component of the corresponding bundle stack.  Use the
four coefficient functors
\[
 R:A\longrightarrow\mathcal D,\qquad
 R':A'\longrightarrow\mathcal D',\qquad
 S:B\longrightarrow\mathcal C,\qquad
 S':B'\longrightarrow\mathcal C'.
\]
The Whittaker twists and characters specifying their targets are
\[
 \begin{array}{c|c|c|c}
 \text{category}&\text{group}&\text{determinant exponent}&
       \text{Whittaker character}\\ \hline
 \mathcal D&G&s&+\\
 \mathcal D'&G&-s&-\\
 \mathcal C&G^\vee&-s'&-\\
 \mathcal C'&G^\vee&s'&+
 \end{array}
\]
in the notation of \eqref{eq:W1}.  The opposite-character pairings
\eqref{eq:W2} give
\[
 \mathcal D'=\mathcal D^\vee,\qquad
 \mathcal C'=\mathcal C^\vee.
\]
Here and below these equalities denote the specified equivalences, and
\(\vee\) denotes categorical duality.

The Ran form of \eqref{eq:L11} identifies
\[
 \mathcal C\simeq\operatorname{KL}_G(a)_{\operatorname{Ran}},\qquad
 \mathcal D'\simeq\operatorname{KL}_{G^\vee}(d)_{\operatorname{Ran}}.
\]
Apply the same comparison to \(-a,-d\), with negative coefficient
characters and positive test characters as at the end of
Section~\ref{sec:ran}.  Its localization exponents differ from the
required opposite exponents by one.  For example,
\[
 -\frac{-a-h^\vee}{2h^\vee}=-s+1.
\]
Tensor by \(L_G^{-1}\) on this localization side and by the opposite
power on its test side; use the corresponding transports for the dual
group.  This gives
\[
 \mathcal C'\simeq\operatorname{KL}_G(-a)_{\operatorname{Ran}},\qquad
 \mathcal D\simeq\operatorname{KL}_{G^\vee}(-d)_{\operatorname{Ran}}.
\]
We use these constructed comparisons throughout.  In particular, no
additional identification between a Kac-module duality and
\eqref{eq:W2} is required.

Let
\[
 L:\mathcal C\longrightarrow A,\qquad
 M':\mathcal D'\longrightarrow B',\qquad
 L':\mathcal C'\longrightarrow A',\qquad
 M:\mathcal D\longrightarrow B
\]
be the four normalized localization functors.  They include the integral
transports just described.  Their normalization is the one used in the
period comparison: induced-stack localization in angle brackets, tensored
with \(l_N^{-1}[\delta_N]\).

\begin{proposition}\label{equivalence:localization}
Each of \(L,M',L',M\) is a continuous localization: its right adjoint is
fully faithful.  Its categorical transpose is therefore fully faithful.
\end{proposition}

\begin{proof}
We use the arbitrary-level localization theorem
\cite[Theorem~13.4.2]{GLCII}.  For a connected reductive group over a field
of characteristic zero, a smooth complete curve, a level \(\kappa\), and
a quasi-compact open \(j:U\hookrightarrow\operatorname{Bun}_H\), this
asserts that
\[
 j^*\operatorname{Loc}_{H,\kappa}:
 \operatorname{KL}(H)_{\kappa,\operatorname{Ran}}
 \longrightarrow D\text{-}\operatorname{mod}_{\kappa}(U)
\]
is a localization.  The target carries the stack twisting associated to
that same level.  Sections~10.1--10.5 of \cite{GLCII} construct this
localization for arbitrary levels; the restriction to the critical level
begins in Section~14 of that work.

Its localization is the Beilinson--Drinfeld localization calculated by
the frame coinvariants in \eqref{eq:T2}.  The arc-equivariant frame
identification \eqref{eq:L1} identifies the affine level with the form
\(-s\operatorname{Kil}\), or its counterpart for the other three
functors.  Proposition~10.1.11 of \cite{GLCII} associates the form
\(-\operatorname{Kil}=2\operatorname{crit}\) with
\(d\log(L_H)\).  Consequently \(-s\operatorname{Kil}\) gives
\(s\,d\log(L_H)\), exactly the determinant twisting displayed here.
Thus the source, the stack twisting, and the localization functor agree
with the ones to which the cited theorem applies.  Tensoring by a
constant invertible line, shifting, and transporting by an integral
power of a determinant line are equivalences and preserve the
localization property.

By the irrational-twist boundedness result of
Section~\ref{sec:foundations}, restriction to a sufficiently large
quasi-compact open is an equivalence for each of \(A,A',B,B'\).  Since
restriction to an open satisfies \(j^*=j^!\), the cited localization
assertion therefore gives the assertion on the full bundle categories.

Finally, for a localization \(Q:\mathcal E\to\mathcal F\),
precomposition identifies functors on \(\mathcal F\) with functors on
\(\mathcal E\) that invert the morphisms inverted by \(Q\).  It is
fully faithful.  Restricting to continuous functors with values in
\(\operatorname{Vect}\) proves that
\(Q^\vee:\mathcal F^\vee\to\mathcal E^\vee\) is fully faithful.
This argument uses the ordinary localization property and does not
require the right adjoint of \(Q\) to preserve colimits.
\end{proof}

\subsection{Construction and the two inverses}

All categories in the following calculation are dualizable.  The period
identity \eqref{eq:T26}, followed by its application to the negated
shifted forms, gives natural isomorphisms
\begin{equation}\label{eq:G1}
 RL\simeq(M')^\vee(S')^\vee,
 \qquad
 R'L'\simeq M^\vee S^\vee.
\end{equation}
These are identities of continuous functors, with the family and
insertion compatibilities proved in Section~\ref{sec:ran}.

\begin{lemma}\label{equivalence:factorization}
There are continuous functors
\(f:A\to(B')^\vee\) and \(f':A'\to B^\vee\), with natural
isomorphisms
\begin{equation}\label{eq:G2}
 \begin{aligned}
 R&\simeq(M')^\vee f,&\qquad (S')^\vee&\simeq fL,\\
 R'&\simeq M^\vee f',& S^\vee&\simeq f'L'.
 \end{aligned}
\end{equation}
They are uniquely determined, together with these isomorphisms, by
\eqref{eq:G1} and the fully faithful transposes of localization.
\end{lemma}

\begin{proof}
Every object of \(A\) is isomorphic to \(L(C)\) for some
\(C\in\mathcal C\).  The first identity of \eqref{eq:G1} therefore
places every value of \(R\) in the essential image of \((M')^\vee\).
Full faithfulness supplies the factorization \(R\simeq(M')^\vee f\),
including its action on morphisms and all coherence data.  The resulting
functor is continuous: for a diagram in \(A\), the comparison from the
colimit of its \(f\)-images to the \(f\)-image of its colimit becomes an
isomorphism after applying \((M')^\vee\), because both that embedding
and \(R\) preserve colimits.  Full faithfulness detects the
isomorphism.  Cancelling the fully faithful embedding in the first
identity of \eqref{eq:G1} gives \((S')^\vee\simeq fL\).
The second identity gives \(f'\) and the other two comparisons by the
same argument, with \(L',R',M^\vee,S^\vee\) in place of
\(L,R,(M')^\vee,(S')^\vee\).
\end{proof}

Let
\[
 e_A:(A')^\vee\xrightarrow{\sim}A,
 \qquad e_B:(B')^\vee\xrightarrow{\sim}B
\]
be the pseudo-identity equivalences of Section~\ref{sec:pseudoidentity},
with shifts \([2\dim Y_G+2\delta_N(G)]\) and
\([2\dim Y_{G^\vee}+2\delta_N(G^\vee)]\), respectively.  They evaluate
a functional on the first slot of the shriek diagonal kernel with the
indicated shift.  The Plancherel identity \eqref{eq:P1} reads
\begin{equation}\label{eq:G3}
 e_A(R')^\vee R\simeq\operatorname{Id}_A,
 \qquad
 e_B(S')^\vee S\simeq\operatorname{Id}_B.
\end{equation}

\begin{proof}[Proof of the equivalence in Theorem~\ref{thm:main}]
Define
\begin{equation}\label{eq:G4}
 \Phi=e_Bf:A\longrightarrow B.
\end{equation}
We display the two inverse candidates to keep their domains and the
role of the two groups explicit:
\[
 \begin{aligned}
 \Psi_{\mathrm{left}}
   &=e_A(f')^\vee M(M')^\vee e_B^{-1}:B\longrightarrow A,\\
 \Psi_{\mathrm{right}}
   &=L(L')^\vee(f')^\vee:B\longrightarrow A.
 \end{aligned}
\]
Transposing \(R'\simeq M^\vee f'\), and using biduality, gives
\((R')^\vee\simeq(f')^\vee M\).  The first identity of
\eqref{eq:G3} and \eqref{eq:G2} consequently imply
\[
 \Psi_{\mathrm{left}}\Phi
 \simeq e_A(f')^\vee M(M')^\vee f
 \simeq\operatorname{Id}_A.
\]
Transposing \(S^\vee\simeq f'L'\) gives
\(S\simeq(L')^\vee(f')^\vee\).  The second identity of
\eqref{eq:G3} therefore implies
\[
 \Phi\Psi_{\mathrm{right}}
 \simeq e_BfL(L')^\vee(f')^\vee
 \simeq\operatorname{Id}_B.
\]
These two inverses agree, since
\[
 \Psi_{\mathrm{left}}
 \simeq\Psi_{\mathrm{left}}\Phi\Psi_{\mathrm{right}}
 \simeq\Psi_{\mathrm{right}}.
\]
Thus \(\Phi\) is fully faithful and essentially surjective.  This uses
neither a composite \(L^RL\) nor a left adjoint on arbitrary
non-holonomic Whittaker tests.

Finally,
\[
 s=\frac{c+h^\vee}{2h^\vee}
   =\frac{c-h^\vee}{2h^\vee}+1,
 \qquad
 -s'=\frac{-1/(rc)-\check h^\vee}{2\check h^\vee}.
\]
Tensoring the source category in Theorem~\ref{thm:main} by the actual
line \(L_G\) identifies it with \(A\), while \(B\) is already the
required target.  Composing this tensor equivalence with \(\Phi\)
proves the theorem in its stated determinant convention.
\end{proof}

\subsection{Localization, coefficients, and families}

\begin{theorem}\label{equivalence:compatibility}
The equivalence \(\Phi\) has the natural localization and coefficient
compatibilities
\begin{equation}\label{eq:G5}
 \begin{gathered}
 \begin{tikzcd}[ampersand replacement=\&,column sep=large,row sep=large]
 \mathcal C \arrow[r,"L"] \arrow[dr,"e_B(S')^\vee"']
   \& A \arrow[d,"\Phi"] \\
   \& B
 \end{tikzcd}
 \\[1ex]
 R\simeq(M')^\vee e_B^{-1}\Phi.
 \end{gathered}
\end{equation}
They hold with the full Ran family structures: they commute with
collisions, insertion of vacua, disjoint products, and the base operations
in the local comparison and period construction.  The first
compatibility uniquely determines \(\Phi\), with its specified
comparison, on the localization quotient.
\end{theorem}

\begin{proof}
Both identities follow from \eqref{eq:G2} and the definition of
\(\Phi\).  Their maps are obtained by lifting the period pairings
through fully faithful embeddings.  Such a lift is unique with its
comparison, so the family compatibilities of \eqref{eq:T26} lift as
well: each coherence diagram becomes the already commuting period
diagram after applying the embedding.  Finally, precomposition with
the localization \(L\) is fully faithful on functor categories, which
proves the uniqueness assertion.
\end{proof}

This gives the unramified localization--Whittaker--Poincar\'e diagram
formulation, as in \cite[Section~4]{GaitsgoryQGL}, with all signs and
shifts fixed.  For any list of distinct insertion points, the
identification \eqref{eq:L11} matches the corresponding KL modules
with the dual-group Whittaker modules.  On that fixed-point category,
the fixed-mark comparison following \eqref{eq:P5} identifies
\(e_B(S')^\vee\) with holonomic shriek Poincar\'e and its continuous
extension.  It is shriek pushforward from the pole model, or from
actual reductions on its strata, with their exponential factors.
The description \eqref{eq:W3} determines the functor on the whole
fixed-point category.  Weyl modules match Whittaker standards with the
label lines of \eqref{eq:L11}; the zero label supplies the vacuum
comparison.

The localization in \eqref{eq:G5} includes
\(l_N^{-1}[\delta_N]\).  With that last normalization removed, its
right-hand side is correspondingly tensored by \(l_N[-\delta_N]\).
The KL comparison uses inverse shifted forms.  After the frame sign and
Cartan shifts in \eqref{eq:L1} and the integral transport above, it
becomes the negative-reciprocal relation between the determinant twists
in Theorem~\ref{thm:main}.  In particular, the compatibility identifies
a functor with specified natural maps, rather than merely identifying
the two categories abstractly.  Its canonicity is relative to the fixed
pinnings, Whittaker characters, orientation normalizations, and theta
characteristic.

\subsection{Components and central characters}

The actual root data give canonical identifications
\[
 X^*(Z_G)\simeq\pi_1(G^\vee),\qquad
 \pi_1(G)\simeq X^*(Z_{G^\vee}).
\]
We use them without passing to a simply connected or adjoint form.
Because the adjoint determinant line has trivial central inertia,
each bundle category decomposes by a central character and a component.

\begin{proposition}\label{equivalence:sectors}
If a source object has central character
\(\eta\in X^*(Z_G)\) and component \(p\in\pi_1(G)\), its image
under \(\Phi\) has central character \(p\) and component \(-\eta\).
Thus, writing character first and component second,
\begin{equation}\label{eq:G6}
 (\eta,p)\longmapsto(p,-\eta).
\end{equation}
\end{proposition}

\begin{proof}
The KL categories over tuples are graded by total central character.
The constant central arcs act using their canonical splitting.  The
localization frame formulas retain this action: changing a frame on
the right by \(h\) acts on its section generators by \(h\), with the
same sign as the infinitesimal formula \eqref{eq:L1}.  Localization
therefore carries the KL central character to bundle central inertia,
without inversion.

In the matched Whittaker cell a label \(\lambda\) gives Cartan bundle
\[
 T_D=T_0\Bigl(-\sum_i\lambda_i x_i\Bigr).
\]
Its component is \(-\sum_i\lambda_i\), since \(2\rho\) lies in the
coroot lattice.  Hence the local comparison sends the total KL central
character to the negative bundle component.  This verification on
cells gives the assertion on all objects: the two decompositions are
already defined on the full Grassmannian, and the open and closed
shriek restrictions to its strata jointly detect the assertion.
Parameter operations preserve both decompositions, so it also holds
on the Ran colimit.

Apply the first identity of \eqref{eq:G5}.  Poincar\'e with the shriek
diagonal preserves bundle components.  A KL input of central character
\(\eta\) therefore gives target component \(-\eta\).  Every source
object of character \(\eta\) is localized from that KL character
sector: localization is essentially surjective and preserves the
character decomposition, so projecting any preimage to the sector
produces a preimage there.  This proves the target component assertion.

For the target character, use \(R\simeq(M')^\vee f\).  The
coefficient functor \(R\) preserves the source bundle component
\(p\).  Under the dual-group local comparison, the part of
\(\mathcal D'\) with that component is the KL character sector
\(-p\).  Since \(M'\) realizes its character sectors, \(f\) sends
a source object of component \(p\) to a functional on \(B'\)
supported on central character \(-p\).  The two factors of the
shriek diagonal kernel have opposite central characters, by its
simultaneous central inertia and the diagonal splitting of the twist.
Evaluation by \(e_B\) consequently gives central character \(p\)
on \(B\), as asserted.
\end{proof}

\subsection{Central torsors and transgression}

We spell out the torsor compatibility because it records the global
form as a groupoid action, including its automorphisms.  Put \(Z=Z_G\).
For a \(Z\)-torsor \(Q\) on \(X\), let \(\tau_Q\) be central
translation on \(Y_G\), with contracted-product convention
\([qz,x]=[q,zx]\).  It canonically preserves the twisting: the
adjoint bundle of the translated bundle is canonically unchanged.
Use the same notation for the corresponding operation for
\(Z_{G^\vee}\)-torsors on \(Y_{G^\vee}\).

Let \(A_{G^\vee}\) be the finite kernel of
\(G^\vee_{\mathrm{sc}}\to G^\vee\).  The root datum gives the
perfect finite pairing
\(Z_G\times A_{G^\vee}\to\mathbb C^\times\).  Use complex
orientation and the exponential with \(2\pi\sqrt{-1}\) to identify
this kernel with the fundamental group.  We normalize the lifting
obstruction so that its integral over \(X\) is the component used in
Proposition~\ref{equivalence:sectors}.  Equivalently, a local
modification with gluing \(t^\lambda\) has relative obstruction
with integral \(-\lambda\).  This is the local lifting obstruction
around an oriented circle, with the same degree convention as the
Cartan bundle in the preceding proof.

The universal bundle on \(X\times Y_{G^\vee}\) has a gerbe of lifts
to \(G^\vee_{\mathrm{sc}}\), with obstruction class
\(\operatorname{obs}\) in degree two and coefficients in
\(A_{G^\vee}\).  Form \([Q]\cup\operatorname{obs}\), with
\(Q\) first, use the finite pairing, and integrate over \(X\).
The resulting class has degree \(1+2-2=1\) and defines a rank-one
local system \(u_Q\) with finite monodromy on \(Y_{G^\vee}\):
\[
 [u_Q]=\int_X\bigl\langle[Q]\cup\operatorname{obs}\bigr\rangle.
\]
This construction is performed on the
torsor and gerbe themselves, as maps to Eilenberg--MacLane spaces,
so it retains isomorphisms, automorphisms, and the coherent tensor law
in \(Q\).  Reversing the two groups gives \(u'_{Q'}\) on \(Y_G\)
for a \(Z_{G^\vee}\)-torsor \(Q'\).  Finite-coefficient comparison
and Riemann existence realize these systems algebraically on charts;
smooth descent gives the systems on the stacks.  Associated-line
monodromy is \(\gamma(z)\) for transport by \(z\) and character
\(\gamma\).  Tensor by these systems means dualizing-normalized
shriek tensor, or ordinary flat tensor inside angle brackets.

\begin{theorem}\label{equivalence:central}
The equivalence \(\Phi\) has coherent natural isomorphisms
\begin{equation}\label{eq:G7}
 \Phi\tau_Q^!\simeq u_Q^{-1}\otimes\Phi,
 \qquad
 \tau_{Q'}^!\Phi\simeq\Phi(u'_{Q'}\otimes-).
\end{equation}
They are natural on the full groupoids of central torsors and respect
their tensor products.  They and \eqref{eq:G6} remain valid after the
integral determinant transport to the convention of
Theorem~\ref{thm:main}.
\end{theorem}

\begin{proof}
We first describe two operations before passing to the global
localization quotient.

\emph{Twisting KL inputs.}
Over a tuple scheme, restrict \(Q\) to the formal neighborhoods of
the marked graphs.  Trivialize it locally and descend KL inputs by the
central arc action.  Denote the resulting operation by \(V\mapsto
V[Q]\).  The same definition works at collisions.  To give precise local
trivializations, write \(D=\sum_i\Gamma_i\subset X\times S\)
for the sum of the marked graphs over the tuple scheme \(S\).
This is finite locally free over \(S\), including at collisions.
The section scheme
\[
 \operatorname{Res}_{D/S}(Q|_D)\longrightarrow S
\]
is \(\acute{e}\)tale and surjective: a geometric fiber admits a
section by choosing a lift on each of its clusters, and formal
\(\acute{e}\)taleness extends this choice across the nilpotents.
This section scheme need not be finite when clusters split.
After passing to it, formal \(\acute{e}\)taleness extends the
universal section uniquely to each thickening \(nD\), and hence
to the completed divisor.  Two such sections differ by a section
of \(Z\) on the completed divisor, locally constant on each
cluster.  Their transitions satisfy the cocycle identity because
they are differences of actual torsor sections.  For insertion or
comparison of tuples, use the sum of their divisors as a common
refinement and restrict its sections.  Unique extension through
thickenings then gives compatibility with the crystal structure,
further collisions, and all resulting coherence diagrams.

On localization charts compare at these chosen lifts of \(Q\).
The out Lie algebra, its differential operators, and the determinant
identifications use the unchanged adjoint bundle.  A change of lift
by \(z\) changes the translated frame by \(z\), precisely the
central transition used in defining \(V[Q]\).  Thus the chart
coinvariant formulas give
\[
 \tau_Q^!L(V)\simeq L(V[Q]).
\]
The comparison holds on bounded-module formulas and their morphisms,
then on their continuous extension; it commutes with insertion and
parameter operations.  The same argument applies to \(M'\) and
\(Q'\) for the dual group.

\emph{Transgression on Whittaker inputs.}
Pull \(u_Q\) to the pole model and tensor there.  Its two pullbacks
along a Whittaker arrow identify canonically.  Indeed the arrow is a
unipotent modification, whose gluing lifts uniquely through the
simply connected cover by the corresponding root subgroups.  Change
the local lift by this lifted modification and retain the lift away
from the marking.  This identifies the lifting gerbes, multiplicatively
under composition of arrows.  Hence tensor by the transgression
system preserves Whittaker equivariance.  Projection formula shows
that coefficient operations and their counits commute with it.

For the local comparison \(\alpha\) of \eqref{eq:L11}, the two
operations are related by
\begin{equation}\label{eq:G8}
 \alpha(V[Q])\simeq u_Q^{-1}\otimes\alpha(V).
\end{equation}
Here the equality includes all tuple-parameter structures.  To check
it, use the Grassmannian model.  Its bundle is identified away from
the markings with \(P_0=\vartheta^{2\rho}\), which has its fixed
lift to the simply connected group.  The lifting obstruction is
therefore relative, supported at the markings.  The chosen disk
sections of \(Q\) trivialize its transgression line: they nullhomotope
\([Q]\) near the support of the relative obstruction, hence nullhomotope
the integrated cup product.  This construction is functorial in the
relative lifting gerbe, so the trivialization holds on the entire
Grassmannian chart, not just on its labeled cells.  On the section
cover, therefore, both sides of \eqref{eq:G8} identify with
\(\alpha(V)\) for every object \(V\), and we use the identity
comparison.  Changing a section by \(z\) on a cluster of total label
\(\lambda\) changes the transgression trivialization by
\((-\lambda)(z)\).

We next check compatibility across collisions.  In an analytic chart of any finite
Grassmannian bound, extend the chosen disk sections to neighborhoods
of the graph unions.  Trivialize the cup product using the relative
obstruction and this section.  Excision makes the trivialization
independent of shrinking the neighborhoods.  Its change is cup
product with the degree-zero change of section followed by
integration, hence evaluation of that section change on the cluster
degree.  That degree is \(-\lambda\) by the obstruction convention.
Neighborhoods and their intersections give these comparisons across
collisions; excision and associativity give the comparisons for
further clustering.  Finite-coefficient comparison transports these
finite-system isomorphisms to the algebraic charts, compatibly with
pullback.

On the KL side the same cluster acts through its total central
character \(\lambda\).  On the Whittaker side \(u_Q^{-1}\) has
transition \(\lambda(z)\), by the calculation just made.  The
component calculation in the proof of
Proposition~\ref{equivalence:sectors} says that \(\alpha\) preserves
these decompositions.  On each finite central-character sector the
transition is a scalar natural transformation, so \(\mathbb C\)-linearity
identifies the two transitions on all objects and morphisms.
For a varying section difference \(z\), its distinct values partition
the completed divisor into open-and-closed pieces; markings with
different values cannot collide.  Use the disjoint-product compatibility
of \(\alpha\) between these pieces, and the single central scalar
on each piece, including all its internal collisions.
The two descent data therefore agree.  The pole--Grassmannian comparison transports the
identification to the Whittaker categories.  It can be checked on
underlying Grassmannian objects because Whittaker equivariance
defines a full subcategory.  This proves \eqref{eq:G8}, with its
coherent parameter compatibilities.  The same construction gives
its dual-group counterpart with \(u'_{Q'}\).

We now pass to the global functor.  Put \(P_B=e_B(S')^\vee\).
Tensoring one slot of the opposite-character pairing by a finite
local system moves that same tensor to its other slot; coefficients
commute with this tensor by the preceding argument.  In the shriek
diagonal kernel, tensor on one factor also moves to the other factor
with the same local system.  Consequently
\[
 P_B(u_Q^{-1}\otimes C)\simeq u_Q^{-1}\otimes P_B(C).
\]
Together with \eqref{eq:G5}, the KL twisting comparison, and
\eqref{eq:G8}, this proves
\(\Phi\tau_Q^!L\simeq u_Q^{-1}\otimes\Phi L\).
Precomposition with the localization \(L\) is fully faithful on
functor categories, so this comparison descends uniquely and gives
the first isomorphism of \eqref{eq:G7}.

For the second, apply the dual-group twisting comparison to
\(R\simeq(M')^\vee f\).  Let \(F\in A\) and \(E\in\mathcal D'\),
viewing \(E\) as a dual-group KL object through the constructed
identification.  Write \(E[Q']\) for the transported central-torsor
action.  The dual-group version of \eqref{eq:G8} gives
\(u'_{Q'}\otimes E\simeq E[(Q')^{-1}]\).  Evaluation on the
localized test \(M'E\) now gives
\[
 \begin{aligned}
 f(u'_{Q'}\otimes F)(M'E)
 &\simeq\bigl\langle R(u'_{Q'}\otimes F),E\bigr\rangle\\
 &\simeq\bigl\langle R(F),u'_{Q'}\otimes E\bigr\rangle\\
 &\simeq\bigl\langle R(F),E[(Q')^{-1}]\bigr\rangle\\
 &\simeq f(F)\bigl((\tau_{Q'}^{-1})^!M'E\bigr).
 \end{aligned}
\]
The first equality is \eqref{eq:G2}; the second is the coefficient
projection formula and the pairing's tensor symmetry; the third is
the dual-group comparison; and the last is the KL localization
compatibility with torsor twisting.  Since \(M'\) is a localization,
these evaluations descend to the natural equality
\[
 f(u'_{Q'}\otimes F)
 \simeq f(F)\circ(\tau_{Q'}^{-1})^!.
\]
To check the variance after applying \(e_B\), let
\(K_B\in B'\otimes B\) be its shifted shriek diagonal kernel,
and put \(\tau=\tau_{Q'}\).  Simultaneous translation preserves
\(K_B\), so
\[
 ((\tau^{-1})^!\otimes\operatorname{Id})(K_B)
   \simeq(\operatorname{Id}\otimes\tau^!)(K_B).
\]
For any continuous functional \(h\) on \(B'\), evaluation of
its first factor consequently gives
\[
 e_B\bigl(h\circ(\tau^{-1})^!\bigr)
   \simeq\tau^!e_B(h).
\]
Taking \(h=f(F)\) proves the second isomorphism.

All maps used in the construction are maps on torsor groupoids:
section descent, the cup product, excision, projection formula, and
diagonal invariance preserve their composition and tensor laws.
The descended comparisons retain these coherences by full
faithfulness.  One can also read the two signs directly on torsor
automorphisms.  An automorphism \(z\in Z_G\) of the trivial
\(Q\) acts on a source object of character \(\eta\) by
\(\eta(z)\).  Its image has component \(-\eta\), so
\(u_Q^{-1}\) has exactly this automorphism.  An automorphism
\(z'\in Z_{G^\vee}\) of the trivial \(Q'\) acts on the
target by \(p(z')\); on the source component \(p\),
\(u'_{Q'}\) has the same automorphism.  Thus the full groupoid
formulas agree with both parts of \eqref{eq:G6}.

Finally, tensoring by \(L_G\) commutes with central
translation, including torsor automorphisms, because this determinant
is formed from the adjoint bundle.  The same is true of central
inertia.  This proves the last assertion and completes the global
unramified correspondence.
\end{proof}
